% Inactive research reconstruction, checkpoint623. Not a book input.
% Scope: the existing chapter's 2 <= m < n, canonical bounded class.
% KN26 already establishes Hermitian Hessian weighted decreasing limits;
% this draft makes no novelty claim and does not certify that whole paper.
\section{Top-degree smooth capacity and weighted positive products}
\label{draft623:top-capacity}

Let $X$ be the chapter's compact Hermitian manifold of dimension $n$,
with $2\leq m<n$. On a coordinate ball $U$ compactly contained in a
holomorphic chart, keep the same Hermitian form $\omega$ and the local
unshifted cone of \cref{hessian:local-unshifted-correction}.
The results of \cref{sec:hessian-smooth-capacity} are inputs at their
stated degrees. The missing torsion estimate at degree $m-1$ can be
proved by summing the diagonal wedge coefficients before taking
absolute values. No positivity of the associated $(n-2,n-2)$-form is
assumed.

Write $S_k(\lambda)$ for the unnormalized elementary symmetric
polynomial and $S_{k;I}$ for the polynomial with coordinates in $I$
omitted. Put $S_0=1$ and $S_k=0$ above the length of the list.

\begin{lemma}\label{draft623:deleted-newton}
If $\lambda\in\overline{\Gamma_m}$, $1\leq q\leq m-1$, and $i\ne j$,
then, with $N=n-2$ and
\[
 c_{N,q}=\frac{q(N-q)}{(q+1)(N-q+1)}<1,
\]
one has
\[
 |S_{q;ij}(\lambda)|^2
 \leq\frac{S_{q;i}(\lambda)S_{q;j}(\lambda)}{1-c_{N,q}}.
\]
This includes $q=N$, when $c_{N,N}=0$.
\end{lemma}
\begin{proof}
First prove the unnormalized Newton inequality for an arbitrary real
list $\mu$ of length $N$. For $1\leq q<N$, put
$R(x)=\prod_{a=1}^N(x+\mu_a)$ and
$G=R^{(N-q-1)}$. Repeated Rolle's theorem makes $G$ a real-rooted
polynomial of degree $d=q+1$. Away from its roots $r_a$, logarithmic
differentiation and Cauchy--Schwarz give
\[
 G'^2-GG''=G^2\sum_{a=1}^d\frac1{(x-r_a)^2}
 \geq \frac1dG^2\left(\sum_{a=1}^d\frac1{x-r_a}\right)^2
 =\frac1dG'^2.
\]
Continuity extends this polynomial inequality to every real $x$.
At zero,
\[
 G=(N-q-1)!S_{q+1}(\mu),\quad
 G'=(N-q)!S_q(\mu),\quad
 G''=(N-q+1)!S_{q-1}(\mu).
\]
Consequently
$S_{q-1}(\mu)S_{q+1}(\mu)\leq c_{N,q}S_q(\mu)^2$.
No sign restriction on $\mu$ or on these coefficients was used.
For $q=N$, its left side is zero and the same assertion holds
with $c_{N,N}=0$.

Apply this to the list with $i,j$ deleted. Put
$A=S_{q-1;ij}$, $B=S_{q;ij}$, $C=S_{q+1;ij}$ and
$D_i=S_{q;i}$, $D_j=S_{q;j}$. Deleted-coordinate positivity gives
$A,D_i,D_j\geq0$: the required degrees are $q+1\leq m$.
Splitting in the two omitted coordinates gives the exact identity
\[
 D_iD_j=A S_{q+1}(\lambda)+B^2-AC.
\]
Since $S_{q+1}(\lambda)\geq0$ and $AC\leq c_{N,q}B^2$, this is
at least $(1-c_{N,q})B^2$. There is no division by $A$ or $D_i$;
all degenerate endpoints are therefore included. Equivalently one
may first add $\epsilon\mathbf1$ and pass to the closed cone.
\end{proof}

\begin{lemma}\label{draft623:top-torsion}
Let $0\leq q\leq m-1$ and let $S$ be a smooth real $(1,1)$-form
in $\overline{\Gamma_m(\omega)}$ pointwise. For smooth real $f$
and nonnegative Borel $a$ supported on a fixed compact coordinate set,
\[
 \left|\int a\,\mathrm df\wedge\dc\omega\wedge S^q\omega^{n-q-2}\right|
 \leq C\left(\int a\,\mathrm df\wedge\dc f\wedge
                           S^q\omega^{n-q-1}\right)^{1/2}
          \left(\int a\,S^q\omega^{n-q}\right)^{1/2}.
\]
The conjugate estimate with $\dc f\wedge\mathrm d\omega$ holds.
For a mixed product, a finite sum of these single-power bounds over
subset sums of its factors suffices.
\end{lemma}
\begin{proof}
For $q=0$ the estimate in \cref{hessian:smooth-torsion} applies.
For $q\geq1$, at one point choose an $\omega$-unitary coframe $e_i$
that diagonalizes $S$. Put $\epsilon_i=\mathrm i e_i\wedge\bar e_i$
and $\Omega=\epsilon_1\cdots\epsilon_n$, so that
$\omega=\sum_i\epsilon_i$ and $S=\sum_i\lambda_i\epsilon_i$.
Writing $B_{ij}=S_{q;ij}(\lambda)$ and $D_i=S_{q;i}(\lambda)$,
ordinary wedge expansion gives
\[
 S^q\omega^{n-q-2}
 =q!(n-q-2)!\sum_{i<j}B_{ij}
                           \prod_{k\notin\{i,j\}}\epsilon_k.
\]
Only the remaining two-plane contributes to its wedge with
$\mathrm df\wedge\dc\omega$. If $\partial f=\sum_i f_i e_i$,
its absolute density relative to $\Omega$ is bounded by
\[
 C\sum_{i<j}(|f_i|+|f_j|)|B_{ij}|.
\]
The constant bounds the tensor coefficients of $\dc\omega$ in all
unitary frames on the fixed compact set. Diagonalization is only
pointwise; no derivative of this moving frame is taken.

The preceding lemma and finite Cauchy--Schwarz bound this sum by
\[
 C\sum_i|f_i|\sqrt{D_i}\sum_j\sqrt{D_j}
 \leq C\left(\sum_i|f_i|^2D_i\right)^{1/2}
          \left(\sum_iD_i\right)^{1/2}.
\]
Here $\sum_iD_i=(n-q)S_q(\lambda)$. Moreover
\[
 \begin{aligned}
 S^q\omega^{n-q}&=q!(n-q)!S_q(\lambda)\Omega,\\
 \mathrm df\wedge\dc f\wedge S^q\omega^{n-q-1}
   &=\kappa q!(n-q-1)!\sum_i|f_i|^2D_i\,\Omega,
 \end{aligned}
\]
where $\kappa=1/(2\pi)$ by the book's $\dc$ normalization.
These two densities are nonnegative. Thus the pointwise torsion
density is bounded by a constant times their geometric mean.
Multiplication by $a$ and integral Cauchy--Schwarz give the assertion,
including the cases where either density vanishes. The other wedge
is conjugate up to a fixed sign, which does not affect its absolute
value. For a mixed product use the exact polarization identity in
\cref{hessian:smooth-torsion}; every subset sum stays in the cone.
No strong positivity at $q=m-1$ is asserted.
\end{proof}

Define $c_m$ and $c_m^*$ by the formulas in
\cref{hessian:smooth-local-capacity}, now with $s=m$.
The smooth CLN bound in \cref{hessian:local-smooth-bounds} makes
this capacity finite on compact subsets. The same suprema show
monotonicity, countable subadditivity and continuity under increasing
unions; the outer-open version agrees on open sets.
The smooth mixed-trace domination in
\eqref{eq:hessian:smooth-capacity-bound} holds at $s=m$, because
it uses only positive mixed traces of at most $m$ cone factors.

\begin{proposition}\label{draft623:top-capacity-convergence}
For the sequence $v_j\downarrow v$ localized as in
\cref{hessian:boundary-localization},
\[
 c_m\{v_k-v>\delta\}\longrightarrow0\qquad(\delta>0).
\]
Every bounded smooth decreasing local sequence has a pointwise limit
quasicontinuous on compact subsets for $c_m$ and $c_m^*$.
Every bounded canonical global admissible potential is locally
quasicontinuous for these capacities.
\end{proposition}
\begin{proof}
Let $q=m-1$, $\theta=\omega^{n-m}$ and $h=v_k-v_j\geq0$ for
$j\geq k$. All differences have support in a fixed $K_0\Subset U$.
The already proved capacity convergence at $s=m-1$, and smooth
mixed-trace domination, give
\[
 \int h\,\ddc g_1\cdots\ddc g_q\omega^{n-q}\longrightarrow0
\]
uniformly in $j\geq k$ and all smooth local cone functions $g_i$
with a fixed common bound. Indeed split at $h=\delta$, bound its
small part by $C\delta$ using CLN, and bound the other part by
$C c_{m-1}\{v_k-v>\delta\}$. Let $k\to\infty$ and then
$\delta\downarrow0$.

Put $w_t=v_j+t h$ and $E_m(w)=\int_U w(\ddc w)^m\theta$.
The fixed smooth collar and the degree-$m$ CLN bound make these
energies uniformly bounded. The exact identity
\eqref{eq:hessian:energy-variation} holds with $s=m$, since its
integration by parts uses only smoothness and compact support of $h$:
\[
 \frac{\mathrm d}{\mathrm dt}E_m(w_t)
 =(m+1)\int h(\ddc w_t)^m\theta
 +m\int h(\ddc w_t)^q
       (\mathrm dw_t\wedge\dc\theta
          -\dc w_t\wedge\mathrm d\theta+w_t\ddc\theta).
\]
In the first-order terms the torsion form is
$(\ddc w_t)^q\omega^{n-q-2}$. The preceding torsion lemma bounds
each integral by the product of a uniformly bounded gradient
factor and a factor tending uniformly to zero by the lower weighted
trace estimate. The gradient bound uses degree $q+1=m$, which is
allowed by \cref{hessian:local-smooth-bounds}. The last term tends
to zero by the absolute mixed-form bound at $q=m-1<m$.
In full,
\[
 \begin{aligned}
 \mathrm d\theta&=(n-m)\mathrm d\omega\wedge\omega^{n-m-1},\\
 \ddc\theta&=(n-m)\ddc\omega\wedge\omega^{n-m-1}\\
 &\quad+(n-m)(n-m-1)\mathrm d\omega\wedge\dc\omega
                                              \wedge\omega^{n-m-2}.
 \end{aligned}
\]
For $n-m=1$ the last term is absent, so no negative power is used.
Thus
\[
 E_m(v_k)-E_m(v_j)=(m+1)\Delta_{kj}+r_{kj},\quad
 \Delta_{kj}=\int_0^1\int h\Lambda_m(w_t)\,\mathrm dt\geq0,
 \quad |r_{kj}|\leq e_k\longrightarrow0.
\]
The inequality $E_m(v_j)\leq E_m(v_k)+e_k$ and a subsequence
approaching the liminf make the bounded energy sequence converge.
Hence $\Delta_{kj}\to0$ uniformly in $j\geq k$.
Since $\Lambda_m(w_t)\geq(1-t)^m\Lambda_m(v_j)$,
\[
 I_m:=\int h\Lambda_m(v_j)\longrightarrow0
                                      \quad\text{uniformly in }j\geq k.
\]

For any smooth competitor $-1\leq a\leq0$, put
\[
 I_\ell=\int h(\ddc v_j)^\ell(\ddc a)^{m-\ell}\theta,\quad
 \eta_\ell=(\ddc v_j)^\ell(\ddc a)^{m-\ell-1},\quad0\leq\ell<m.
\]
These integrals are nonnegative. Exact first-order Stokes gives
\[
 \begin{aligned}
 I_\ell&=-\int\mathrm dh\wedge\dc a\wedge\eta_\ell\theta
             +\int h\,\dc a\wedge\eta_\ell\mathrm d\theta,\\
 F_\ell(h)&:=\int\mathrm dh\wedge\dc h\wedge\eta_\ell\theta\\
 &=I_{\ell+1}-\int h\ddc v_k\wedge\eta_\ell\theta
                +\int h\,\dc h\wedge\eta_\ell\mathrm d\theta.
 \end{aligned}
\]
The second integral in the first line is $o(1)$ uniformly in $a$ and
$j\geq k$. Apply mixed polarization in the torsion lemma: subset
sums of the factors in $\eta_\ell$ have one common local bound,
their gradient factors are bounded by CLN and their lower traces
with weight $h$ tend to zero. The last integral has the same
property. For that term the first factor is bounded by the energy
of $h$, at most twice the sum of those of $v_k,v_j$ for each
positive single-power gradient pairing; the second uses
$h^2\leq C h$ and the lower weighted trace estimate.

The form $\eta_\ell\theta$ has positive gradient pairing: appending
one nonnegative rank-one form completes a polarized product of
$m$ cone forms. Consequently Cauchy--Schwarz for this positive
$(n-1,n-1)$-form and the uniform gradient bound for $a$ give
\[
 I_\ell\leq C F_\ell(h)^{1/2}+o(1),\qquad
 F_\ell(h)\leq I_{\ell+1}+o(1).
\]
The dropped integral is nonnegative. Starting with $I_m\to0$,
finite backward induction gives $I_0\to0$ uniformly over all
competitors. Chebyshev's inequality and continuity of $c_m$ under
increasing unions yield
\[
 c_m\{v_k-v>\delta\}
   =\lim_{j\to\infty}c_m\{v_k-v_j>\delta\}
   \leq\delta^{-1}\sup_{j\geq k,a}I_0\longrightarrow0.
\]

Choose a subsequence $k_N$ with
$c_m(G_N)<\epsilon2^{-N}$ for the open sets
$G_N=\{v_{k_N}-v>2^{-N}\}$. These sets are open because the
pointwise decreasing limit $v$ is upper semicontinuous.
Outside their union, $v_{k_N}\to v$ uniformly, proving continuity
there. Countable subadditivity bounds the exceptional capacity by
$\epsilon$. Boundary localization gives this conclusion on each
compact subset of the original local ball; open-set equality also
gives $c_m^*$. For a bounded global admissible $u$, apply its global
smooth decreasing approximation and the local correction $u+\rho$
from \cref{hessian:local-unshifted-correction}. Subtracting smooth
$\rho$ preserves continuity on the complement of the exceptional
open set and identifies the limit with the canonical values of $u$.
\end{proof}

\begin{lemma}\label{draft623:positive-capacity-domination}
On $K\Subset U$, the positive shifted currents
$Q_p(u_1,\ldots,u_p)$, $1\leq p\leq m$, for globally bounded
admissible potentials have trace and coefficient variations bounded
on Borel $E\subset K$ by $C c_m^*(E)$. The constants are uniform
when all potentials have a common bound.
\end{lemma}
\begin{proof}
First take smooth potentials, put $\sigma_i=\omega+\ddc u_i$, and
choose the fixed local correction $\tau=\ddc\rho\geq2\omega$.
Then
$\beta_i=\ddc(u_i+\rho)=\sigma_i+(\tau-\omega)$ lies in the
same local cone. Since $\tau-\omega\geq0$, mixed positivity gives
\[
 0\leq\int_E Q_p\wedge\omega^{m-p}
    \leq\int_E\beta_1\cdots\beta_p\omega^{n-p}
    \leq C c_m(E),\qquad E\subset K.
\]
Here $Q_p\wedge\omega^{m-p}=\sigma_1\cdots\sigma_p\omega^{n-p}$.
The last bound follows from smooth mixed-trace domination at $s=m$
for the bounded corrected potentials and a normalized bounded sum with a fixed strictly
plurisubharmonic function. Positivity of $Q_p$ bounds each of its
coefficient variations on the compact coordinate set by a fixed
multiple of its trace measure.

For bounded potentials choose the smooth decreasing global
approximations already established in the chapter. The joint
convergence in \cref{hessian:joint-decreasing-products} identifies
their positive-current limit with the previously constructed $Q_p$.
On a compact neighborhood of $K$ inside $U$, lower semicontinuity
on open sets preserves trace and coefficient variation estimates.
For an arbitrary open $G\supset E$ intersect with a fixed relatively
compact neighborhood of $K$ before taking this limit; its capacity
is at most $c_m(G)$. Therefore every such trace or variation measure
$\mu$ satisfies $\mu(E)\leq Cc_m(G)$. Infimize over $G$ to obtain
the asserted outer-open bound. The constants of the smooth estimates
are independent of the approximation index. The argument does not
multiply a rough current of positive order or use an unproved rough
diagonal polynomial expansion.
\end{proof}

\begin{theorem}\label{draft623:weighted-positive-products}
Suppose $a_j\downarrow a$ and $u_i^j\downarrow u_i$ are bounded
canonical global admissible potentials with common bounds. Then,
for $1\leq p\leq m$,
\[
 a_j Q_p(u_1^j,\ldots,u_p^j)
        \longrightarrow a Q_p(u_1,\ldots,u_p)
\]
weakly as currents with measure coefficients. At $p=m$ this is weak
convergence of weighted Radon measures. No order-zero assertion for
the unwedged $D_m$ is made.
\end{theorem}
\begin{proof}
Work with one coefficient measure $\nu_j\to\nu$ on a coordinate
chart; its weak convergence and local mass bound follow from
\cref{hessian:joint-decreasing-products}. The preceding lemma bounds
all these variations by $C c_m^*$ on a fixed compact test set.
The capacity proposition makes every canonical $a_j$ and $a$
quasicontinuous for this capacity. Allocate exceptional capacities
summably to this countable family and take their open union $G$
with $c_m(G)<\epsilon$. On the compact $F=K\setminus G$, all the
weights are continuous and Dini's theorem gives uniform convergence
$a_j\to a$. Extend $a|_F$ to a continuous $g$ on $U$ with the same
bound $M$ by Tietze's theorem, taking $g=0$ if $F$ is empty.
For a fixed $\phi\in C_c(U)$ with support in $K$, subtraction gives
\[
 \begin{split}
 \left|\int\phi(a_j\,\mathrm d\nu_j-a\,\mathrm d\nu)\right|
 &\leq\|\phi\|_\infty\|a_j-a\|_F|\nu_j|(K)
      +\left|\int\phi g\,\mathrm d(\nu_j-\nu)\right|\\
 &\quad+2M\|\phi\|_\infty
                \bigl(|\nu_j|(K\cap G)+|\nu|(K\cap G)\bigr).
 \end{split}
\]
The first two terms tend to zero for fixed $G$, and the last is at
most $4MC\|\phi\|_\infty\epsilon$. Let $\epsilon\downarrow0$.
Continuous extension is tested against the whole measure, avoiding
any weak-limit claim for a discontinuous closed-set indicator.
Repeat coefficientwise and use a finite partition of unity.
The positive top-degree current is already a Radon measure, so the
same argument gives the claimed weighted convergence at $p=m$.
\end{proof}
